\parttitle{V}{Hardware Considerations}

\noindent
This part reads the platform landscape through one question: \emph{which
hardware properties bind for reservoir computing specifically?} Reservoirs do
not need deep coherent circuits, error-corrected logic, or trained on-device
parameters; they need fast repeated state preparation and measurement, good
mid-circuit measurement and reset where recurrence is wanted, enough shots to
beat Eq.~\eqref{eq:shot-noise}, and noise of a \emph{kind} that contracts
without erasing (Section~\ref{sec:channels}). Those priorities rank platforms
differently from, say, fault-tolerance roadmaps, and the ranking is
task-dependent. Quantitative gate-fidelity and coherence figures below are
order-of-magnitude characterisations of current practice; exact values change
with every calibration cycle and should be read from the provider's live
calibration data at run time, as the Part~IV pipeline does when pointed at a
real backend.

\section{IBM superconducting processors (primary target)}
\label{sec:hw-ibm}

IBM's fixed-frequency transmon processors are this document's primary
hardware target because they combine the largest publicly accessible
register counts, the most mature dynamic-circuit support, and the toolchain
(Part~IV) whose abstractions map one-to-one onto the physics. The
QRC-relevant profile:

\emph{Gate set and topology.} Native alphabet \code{ecr}/\code{rz}/\code{sx}/\code{x}
on a heavy-hex lattice: each qubit couples to two or three neighbours. For
reservoirs this makes brickwork-local circuits (Listings~\ref{lst:cirq} and
\ref{lst:qiskit}) essentially free to route, while any nonlocal coupling
costs SWAP chains --- the depth-78/32-\code{ecr} accounting of
Section~\ref{sec:transpilation} is the honest unit of thinking. Single-qubit
error rates sit around $10^{-4}$, two-qubit around a few $10^{-3}$; the
two-qubit count therefore dominates any reservoir step's error budget.

\emph{Coherence and clocks.} Transmon $T_1, T_2$ of order
$10^2\,\mu$s against $\mu$s-scale gate layers permit circuit
durations of hundreds of layers before decoherence saturates --- but the
reservoir-specific point is that NISQRC-style operation
(Section~\ref{sec:midcircuit}) makes usable \emph{memory} independent of
this budget, which was demonstrated on exactly this platform
\citep{hu2024overcoming}. Wall-clock throughput --- circuits
$\times$ shots per second --- is the binding resource for time-series work,
and superconducting platforms lead it by orders of magnitude.

\emph{Mid-circuit measurement and reset.} Production capability with
per-step classical registers and classically conditioned gates
(\code{if\_else}); measurement durations of order $\mu$s with
active reset. This is the primitive that separates genuinely recurrent
hardware reservoirs from windowed feature maps, and its error (read-out
confusion of order $10^{-2}$, plus measurement-induced dephasing of
spectators) should be characterised per-qubit before trusting long
sequences; the natural-reservoir demonstrations of
\citet{suzuki2022natural,yasuda2023quantum} and the NISQRC experiment
\citep{hu2024overcoming} are the platform's QRC track record, alongside the
utility-scale dynamics simulations that certify large-register control
\citep{kim2023evidence}.

\emph{Error mitigation.} Estimator-level mitigation (read-out
mitigation, zero-noise extrapolation, probabilistic error cancellation)
applies directly to feature estimation, at shot-multiplier costs that must be
budgeted like any other shots~\citep{larose2022mitiq}. For reservoirs a
cheaper posture often suffices: keep circuits shallow, measure local
observables, and let the read-out's regularisation absorb residual bias ---
mitigation is worth its multiplier mainly when features must be compared
across sessions.

\emph{Access economics.} Queued cloud access with per-second pricing on
premium plans; batched jobs (Part~IV's shape) amortise queue overhead.
Session/batch primitives keep calibration approximately fixed within an
experiment --- important because a noise-defined reservoir drifts between
calibrations (Section~\ref{sec:arch-dissipative}).

\section{Quantinuum trapped-ion systems}
\label{sec:hw-quantinuum}

The H-series QCCD architecture transports ions between gate zones, giving
effective all-to-all connectivity, two-qubit fidelities at the
$99.8$--$99.9\%$ level, and mid-circuit measurement with conditional logic as
a routine primitive \citep{moses2023race}. For QRC: the fully connected
disordered Hamiltonian of Eq.~\eqref{eq:fn-hamiltonian} compiles without
routing overhead, and measurement/reset quality is the best available ---
the natural platform for \emph{clean} recurrent-reservoir physics
experiments (dynamical-regime studies, ESP characterisation) where per-step
fidelity matters more than throughput. The cost is clock rate: transport and
cooling make shot rates orders of magnitude below superconducting, so
long-time-series benchmarks are expensive; restart-protocol studies with
short windows, or physics studies with modest $T$, fit best. No published
QRC benchmark at the scale of the superconducting or neutral-atom
demonstrations exists as of this writing --- flagged as an opportunity in
Part~VIII.

\section{IonQ trapped-ion systems}
\label{sec:hw-ionq}

Photonic-interconnect-oriented ion traps with all-to-all gates within a
chain, accessed via native API and Braket. Similar QRC reasoning to
Section~\ref{sec:hw-quantinuum} with smaller registers and, historically,
less exposed mid-circuit control; algorithmic-qubit marketing figures should
be translated back into raw two-qubit fidelity and shot throughput before
budgeting an experiment.

\section{Rigetti superconducting processors}
\label{sec:hw-rigetti}

Tunable-coupler transmons with faster two-qubit gates but historically lower
fidelities than IBM; native parametric access and low-latency classical
integration make it a reasonable testbed for feedback reservoirs
(Section~\ref{sec:arch-measurement}) where round-trip latency to classical
logic is the experiment. For the pipelines of Part~VII the IBM stack's
maturity currently wins on reproducibility.

\section{Neutral-atom analog platforms (QuEra, Pasqal)}
\label{sec:hw-neutral}

Rydberg tweezer arrays (Section~\ref{sec:arch-rydberg}) invert several
trade-offs: hundreds of qubits with programmable \emph{geometry}, analog
evolution with no per-gate compilation error, but destructive fluorescence
read-out --- every time step is a full re-preparation, and there is no
mid-circuit measurement in the digital sense. The 108-qubit QRC demonstration
of \citet{kornjaca2024large} is the largest to date and defines the
platform's role: large feature banks via ensemble/restart operation on
classification-style or short-window temporal tasks, with geometry and
detunings as reservoir hyperparameters, and logical-array progress
\citep{bluvstein2024logical} indicating the control ceiling is still rising.
Per-step latency (array reload of order tens of ms) makes long
autoregressive forecasting loops the platform's weak axis.

\section{Photonic platforms (Xanadu-class CV)}
\label{sec:hw-photonic}

Time-multiplexed squeezed-light processors \citep{madsen2022quantum} and
tabletop multimode systems (Section~\ref{sec:arch-cv}) offer
room-temperature operation, natural weak/continuous measurement, and clock
rates far beyond matter qubits --- attractive for streaming-signal reservoirs
--- with loss as the dominant decoherence and non-Gaussian resources as the
scarce ingredient. Experimental optical QRC with controllable memory
\citep{paparelle2025experimental,garciabeni2023scalable} and large-scale
Gaussian-boson-sampler feature maps \citep{cimini2025gaussian} are the
current state of practice; microwave-domain analog circuit QED
\citep{senanian2024microwave,dudas2023quantum} plays the same role for RF
signals.

\section{Cross-platform summary}

\begin{table*}[h]
\centering\small
\caption{Platform comparison on QRC-binding axes (orders of magnitude,
mid-2026 practice). ``MCM'' = mid-circuit measurement and reset quality;
``throughput'' = circuits$\times$shots per wall-clock unit.}
\label{tab:hw-compare}
\begin{tabular}{@{}p{2.9cm}p{2.2cm}p{2.0cm}p{2.0cm}p{2.2cm}p{2.6cm}@{}}
\toprule
Platform & Connectivity & MCM & Throughput & Noise character & QRC record \\
\midrule
IBM superconducting & heavy-hex, local & production & highest & depolarising $+$ read-out; $T_1$ decay & NISQRC; natural reservoirs \citep{hu2024overcoming,suzuki2022natural,yasuda2023quantum} \\
Quantinuum ion & all-to-all & excellent & low & low, well-characterised & (open) \\
IonQ ion & all-to-all (chain) & limited exposure & low & low & early studies \\
Rigetti supercond. & lattice, local & available & high & higher gate error & feedback testbed \\
QuEra/Pasqal neutral & geometry-program. & none (destructive) & medium (batch) & analog drift; re-prep. & 108-qubit QRC \citep{kornjaca2024large} \\
Photonic CV & mode-mesh & continuous/weak & very high & loss & optical QRC \citep{paparelle2025experimental,cimini2025gaussian} \\
\bottomrule
\end{tabular}
\end{table*}

Two decision rules compress the table. If the task is a \emph{long temporal
sequence} (forecasting loops, streaming), throughput and MCM dominate:
superconducting first, photonic for native-signal front-ends. If the task is
a \emph{physics question about reservoirs} (dynamical regimes, ESP,
coherence's role), fidelity and controllability dominate: trapped ions and
neutral atoms first, with the analog platforms unique for
geometry-as-hyperparameter questions. In every case the noise model is part
of the reservoir: re-characterise it per session, and treat cross-session
comparisons as different reservoirs unless proven otherwise.
